% Compile with python build.py. The build stages the vector figures in plates/.
\documentclass[11pt,a4paper]{article}
\usepackage[inner=25mm,outer=25mm,top=23mm,bottom=23mm,headheight=13pt,headsep=6mm]{geometry}
\usepackage{lmodern,amsmath,amssymb,graphicx,booktabs,array,xcolor,caption,fancyhdr}
\usepackage{microtype}
\usepackage{hyperref}
\graphicspath{{plates/}}
\newcommand{\plate}[2]{\includegraphics[width=#2\linewidth]{plates/#1.pdf}}
\definecolor{ink}{HTML}{191B1B}
\definecolor{quiet}{HTML}{515555}
\definecolor{rulegray}{HTML}{A5AAA6}
\definecolor{linkblue}{HTML}{233C83}
\hypersetup{pdftitle={How a puzzle grows},pdfauthor={Formverk Studio research notes},
 pdfsubject={An illustrated technical companion to the Nervous System puzzle paper},
 colorlinks=true,linkcolor=ink,citecolor=linkblue,urlcolor=linkblue,bookmarksnumbered=true}
\captionsetup{font=small,labelfont=bf,labelsep=period,justification=raggedright,singlelinecheck=false,skip=6pt}
\setlength{\parindent}{0pt}
\setlength{\parskip}{4pt plus 1pt minus 1pt}
\setlength{\textfloatsep}{9pt plus 2pt minus 2pt}
\setlength{\intextsep}{10pt plus 2pt minus 2pt}
\renewcommand{\topfraction}{.98}
\renewcommand{\bottomfraction}{.95}
\renewcommand{\textfraction}{.03}
\pagestyle{fancy}\fancyhf{}
\fancyhead[L]{\small\scshape Formverk Studio\enspace/\enspace Research notes}
\fancyhead[R]{\small\itshape Generative jigsaw geometry}
\fancyfoot[L]{\footnotesize\color{quiet}How a puzzle grows}
\fancyfoot[R]{\small\thepage}
\renewcommand{\headrulewidth}{.35pt}
\newcommand{\rubric}[1]{{\small\sffamily\bfseries\color{quiet}#1}}
\newcolumntype{L}[1]{>{\raggedright\arraybackslash}p{#1}}
\raggedbottom
\begin{document}
\rubric{ILLUSTRATED TECHNICAL COMPANION}\hfill{\small 28 November 2026}
\par\vspace{3pt}
{\fontsize{32}{34}\selectfont\bfseries How a puzzle grows\par}
{\Large Shared boundaries in generative wooden jigsaws\par}
\vspace{2pt}
{\small A reading of the 2012 model by Jesse Louis-Rosenberg, Alec Resnick and
Jessica Rosenkrantz, with a comparison to the independent Formverk implementation.}
\par\vspace{3pt}\hrule\vspace{4pt}
\textbf{Abstract.}
One cut through a sheet of plywood makes two complementary edges. The design
problem is to choose their shared curve so that neighboring pieces interlock
and narrow wooden features survive fabrication. We follow a photograph and a
cut plan through production, examine the phase and temperature fields of the
published Dendrite model, and compare six constructions of a single seam.
A final implementation study traces how Formverk recovers, exports and checks
the geometry. The diagrams distinguish published mechanisms, local measurements
and illustrative constructions.
\par\smallskip
{\small\textit{Keywords}\enspace Phase-field model, elastic rods, planar partition,
laser cutting, computational fabrication.}
\section{Two paths from image to object}
\label{sec:overview}
A customer supplies artwork, chooses a size, outline and cut style, and can
place recognizable figure pieces called \emph{whimsies}~\cite{configurator}.
In its original series the studio mounted a print on plywood. A later series
used direct UV printing on maple~\cite{puzzleproject,wave}. The chosen board
and fixed figures constrain the construction of a planar partition. At the
laser, the paths meet physically (Fig.~\ref{fig:route}).\par
\begin{figure}[!ht]\centering
\plate{production-route}{1}
\caption{Material and geometry in the studio's documented workflow.
The upper branch carries the artwork on wood. The lower branch carries a board
boundary and internal cut paths. Arrows identify what passes between stages.}
\label{fig:route}\end{figure}\par
\textbf{Evidence and scope.} The primary mathematical source is the authors'
technical draft~\cite{puzzle2012}. The studio also publishes a fabrication
account and descriptions of all six cuts~\cite{puzzleproject,cutguide}.
The deployed service layout, current parameter sets and dimensional clearance
thresholds remain undocumented in these sources. Formverk is examined separately
through its own source code~\cite{formverk}.

\clearpage
\section{Dendritic solidification}
\label{sec:dendrite}

The 2012 Dendrite model starts with the whole puzzle already divided into
cells. Seeds expand by reaction and diffusion until neighboring regions meet.
A fixed figure reflects the diffusing fields, so a cell can wrap around its
outline. Two sinusoidal perturbations, one coarse and one fine, ripple the
initial seams. Their wavelengths and amplitudes help determine where branches
form \cite{puzzle2012}.

\paragraph{Why a ripple grows.}
Let \(p(\mathbf{x},t)\) be a local phase value between zero and one and
\(T(\mathbf{x},t)\) the dimensionless temperature at position \(\mathbf{x}\)
and time \(t\). The authors print the isotropic two-phase model in the form
\cite[Eqs.~5--7]{puzzle2012}
\begin{equation}
\begin{aligned}
  \tau\,\partial_t p
    &= \epsilon^2\nabla^2p+p(1-p)\bigl(p-\tfrac12+m(T)\bigr),\\
  \partial_t T
    &= \nabla^2 T+K\,\partial_t p,
    & m(T)&=\frac{\alpha}{\pi}\arctan\!\bigl(\gamma(T-T_e)\bigr).
\end{aligned}
\label{eq:dendrite-thermal}
\end{equation}
Here \(\epsilon\) sets the width of the diffuse interface, \(\tau\) its
relaxation time, and \(K\) the latent-heat coupling. The constants
\(\alpha\) and \(\gamma\) set the strength and temperature sensitivity of the
bias \(m\). The authors use \(T_e=0\) for equilibrium. At a growing tip,
heat diffuses sideways faster than the seam advances. The flanks warm
while the tip continues into colder material, amplifying the initial ripple
(Fig.~\ref{fig:thermal}).

\begin{figure}[!ht]\centering
\plate{thermal-feedback}{0.76}
\caption{Thermal feedback at an advancing tip. A protrusion reaches colder
material ahead. Heat from solidification spreads sideways and slows growth on
the flanks. The signs mark reinforcement and inhibition.}
\label{fig:thermal}
\end{figure}

\textit{Reading the sign convention.} The draft labels $p=0$ solid and
$p=1$ liquid, yet its positive $K\,\partial_t p$ term supplies heat as $p$
increases. Its Fig.~2 describes heat release on solidification. Equation~\eqref{eq:dendrite-thermal}
preserves the printed signs. A numerical reconstruction must choose a
consistent phase and heat convention~\cite[\S3.1]{puzzle2012}.

\paragraph{How neighboring pieces share a seam.}
A single solid--liquid front has only two sides. At a junction of three
pieces, each region needs a distinct identity. The model assigns each
active label \(i\) a phase fraction \(p_i\) and its own temperature \(T_i\),
with \(\sum_i p_i=1\). The draft combines interactions between pairs
\cite[Eqs.~8--10]{puzzle2012}. In compact notation,
\begin{equation}
  m_{ij}=\tfrac12\bigl[m(T_i)+m(T_j)\bigr],
  \qquad
  R_i=\sum_{j\ne i}p_i p_j\bigl(p_i-p_j+m_{ij}\bigr).
\label{eq:dendrite-pairwise}
\end{equation}
Here \(R_i\) denotes the summed local reaction for label \(i\). The factor
\(p_i p_j\) localizes each term to places where two phases overlap. With the
printed symmetric bias, the two ordered reactions for a pair sum to
$2p_i p_j m_{ij}$. An implementation must therefore also account for the
stated constraint $\sum_i p_i=1$. For
production-scale layouts the authors reused five labels across nonadjacent
pieces. Heat can then travel between pieces sharing a label and soften fine
branches in long runs. After growth, the studio traced the phase grid into
cut paths and smoothed tight folds to protect narrow wooden necks
\cite{puzzle2012,puzzleproject}.

\clearpage
\section{Six ways to construct a shared seam}
\label{sec:seams}

A cut between two pieces gives each the outline of its neighbor. The six styles
in the custom maker differ in how they construct that line. For the first three,
the starting regions already fill the board before the characteristic interlocks
develop. Figure~\ref{fig:growing-seams} follows a single edge through that change.

\begin{figure}[!ht]
\centering
\begin{minipage}{.98\linewidth}
\textbf{(a) Dendrite.}
Seeds spread until their territories meet. Small waves disturb the new boundary.
The 2012 solver gives adjoining pieces different phase labels and reuses five
labels across the board. Each phase carries a temperature field. As a tip advances,
it releases heat that slows growth along its flanks. That feedback lets branches
develop on either side of the seam. Over long runs, heat can spread between
distant pieces sharing a label and soften the branches~\cite{puzzle2012}.
\par\smallskip
{\centering\plate{dendrite}{.89}\par}
\medskip
\textbf{(b) Amoeba.}
The studio uses the same crystal-growth model with a different parameter set.
The paper identifies seed spacing, ripple wavelength and amplitude, and
supercooling as controls over the resulting shapes. The published description
of Amoeba describes short, closely spaced projections and a higher piece density.
Its numerical preset remains unpublished~\cite{puzzle2012,cutguide}.
\par\smallskip
{\centering\plate{amoeba}{.89}\par}
\medskip
\textbf{(c) Maze.}
Seed regions first expand around the outline and fixed figures. Their shared
edges become elastic rods. Each rod tends to straighten, but its length increases
inside a confined area. The edges buckle, collide and push one another into the
long folds visible in the finished wood~\cite{maze}.
\par\smallskip
{\centering\plate{maze}{.89}\par}
\end{minipage}
\caption{Three routes from initial territory to an interlocking edge. The first two
rows share the studio's phase-field mechanism. The third illustrates the excess
length accommodated by a growing rod. The drawings are conceptual.}
\label{fig:growing-seams}
\end{figure}

\clearpage
\subsection{Curling a rod, joining arcs and trading cells}

Wave adds a force to the elastic edges of Maze. Coral and Matrix instead compare
candidate arrangements and score the interlock they create. The unit of change
in Figure~\ref{fig:constructed-seams} is, respectively, a growing rod, an arc
chain or a grid cell.

\begin{figure}[!ht]
\centering
\begin{minipage}{.98\linewidth}
\textbf{(a) Wave.}
A directional Voronoi layout elongates the initial pieces along a flow field.
Each edge receives vortices according to its length. They move with the edge as
the rod lengthens, drawing it into the spirals and hammerhead-shaped hooks
characteristic of this cut~\cite{wave,cutguide}.
\par\smallskip
{\centering\plate{wave}{.89}\par}
\medskip
\textbf{(b) Coral.}
Jules Khan's generator starts with neighboring cells and proposes replacement
edges made of joined circular arcs. Random walks yield several routes between
the same endpoints. The search scores arc chains for length and interlock,
avoids self-crossings and leaves enough material in narrow parts of a piece
for it to hold together~\cite{coral,cutguide}.
\par\smallskip
{\centering\plate{coral}{.89}\par}
\medskip
\textbf{(c) Matrix.}
The outline, fixed figures and guide curves orient a mesh of mostly quadrilateral
cells. The initial pieces each take roughly the same number of cells. The
optimizer looks several cell swaps ahead as it balances size and interlock. To
score a shared edge, it maps its tangent directions onto a circle. The studio
treats coverage beyond \(180^\circ\) as evidence of interlock. A range approaching
\(360^\circ\) indicates a stronger grip~\cite{matrix,cutguide}.
\par\smallskip
{\centering\plate{matrix}{.89}\par}
\end{minipage}
\caption{Three constructions at edge scale. The red marks in (a) represent
vortices. In (b), gray arcs trace alternative routes and blue marks the selected
seam. In (c), the highlighted cells change ownership across the boundary.}
\label{fig:constructed-seams}
\end{figure}

\clearpage
\section{From photograph to plywood}
\label{sec:fabrication}

The photograph supplies the printed image and can suggest the perimeter.
Inside that perimeter, fixed figures reserve space before the surrounding
pieces grow. Figure~\ref{fig:production-studies} follows the geometry through
three decisions that affect the physical puzzle.

\begin{figure}[!ht]
\centering
\begin{minipage}{.98\linewidth}
\textbf{(a) Find an outline.}
Nervous System credits rembg for its Magic Shape feature~\cite{magicshape}.
In the archived replay, U\({}^2\)-Net takes a \(320\times320\) RGB image, fuses
predictions at six scales and returns a foreground mask~\cite{u2net}. The red
line below traces the largest outer contour at a mask value of 128. On this
one photograph, the full and compact models matched the upstream reference
mask with intersection-over-union scores of 0.99916 and 0.97213. The studio
has not published how it converts its prediction into a cut perimeter.
\par\smallskip
{\centering\includegraphics[width=.91\linewidth]{mask-flow.png}\par}
\medskip
\textbf{(b) Reserve the figure pieces.}
For \emph{Herding Cats}, the studio used Deepnest to fit cat silhouettes inside
a larger cat outline, then adjusted their placement by hand~\cite{herdingcats}.
The surrounding pieces grow through the gaps. In the published Dendrite
method, fixed figures act as reflecting boundaries during initialization~\cite{puzzle2012}.
\par\smallskip
{\centering\plate{fixed-figures}{.89}\par}
\medskip
\textbf{(c) Make a cuttable boundary.}
The original Dendrite workflow dilates pieces and removes fragments that pinch
off. Modified marching squares traces the resulting contours. Where a point
lies too close to another edge, local smoothing widens the clearance before
export~\cite{puzzleproject}. The laser follows those paths across the printed
wood. The studio also developed an air-assist nozzle to limit scorching while
cutting its puzzles~\cite{airassist}.
\par\smallskip
{\centering\plate{cut-to-wood}{.78}\par}
\end{minipage}
\caption{The photo, mask and measured contour in (a) use the archived
\href{https://github.com/danielgatis/rembg/blob/332c16b58e0194f0960f9d6cf38746e7be8546b7/examples/animal-1.jpg}{rembg example photograph}.
The schematics show cells meeting a fixed silhouette in (b) and one cut
producing two matching wooden edges in (c).}
\label{fig:production-studies}
\end{figure}

\clearpage
\section{An independent implementation}
\label{sec:local}
Formverk Studio and the command line call the same puzzle generator.
The request specifies an outline, dimensions in millimetres, a target count,
a random seed and any fixed figures. The selected style determines the geometry
that the exporter receives (Fig.~\ref{fig:engine})~\cite{formverk}.\par
\begin{figure}[!ht]\centering
\plate{local-engine}{1}
\caption{Eight local cut styles use four solver families. All return polygon
faces and shared cut paths to one exporter. Reopening the written SVG tests
the geometry that downstream readers receive.}
\label{fig:engine}\end{figure}\par
\paragraph{How the four families construct a seam.}
Dendrite, Amoeba, Linear and Grid grow a raster partition. After phase and
temperature evolve, the solver repairs labels and traces their shared boundaries.
Maze and Wave grow elastic rods along shared edges. Wave also applies a curl
field. The studio credited PositionBasedDynamics for its rod simulations in
2017~\cite{maze,rods2016}. Formverk's Coral searches for chains of tangent
circular arcs that clear nearby cuts. Matrix grows connected groups of mesh
cells, then exchanges boundary cells between neighbors to improve interlock.
\paragraph{One output boundary.}
All families share the millimetre-scale cut and artwork exporter. It reloads
\texttt{cut.svg}, rebuilds faces, and records coverage, duplicate segments,
kerf and neck checks in \texttt{audit.json}. It estimates interlock by sampling
360 translation directions in the plane. The player receives outlines and
neighbors in \texttt{puzzle.json}. These tests measure planar geometry.
Physical holding strength requires measurements on fabricated pieces.\par
\begin{table}[!ht]\centering\small
\caption{Evidence for three checks. IoU denotes intersection over union of
thresholded foreground sets. Each conclusion is limited to the stated specimen.}
\label{tab:checks}
\begin{tabular}{@{}L{.21\linewidth}@{\hspace{.025\linewidth}}L{.35\linewidth}@{\hspace{.025\linewidth}}L{.36\linewidth}@{}}
\toprule
\textbf{Check} & \textbf{Data and observation} & \textbf{Interpretation}\\
\midrule
Single crystal & $40\times40$ cells, ten time steps. All cells passed the
upstream comparison. & Reproduces the independent FiPy reference case~\cite{fipy}.\\[3pt]
Photo mask & One upstream example, threshold 128. Full-model IoU 0.99916,
compact-model IoU 0.97213. & Measures agreement with the distributed reference
mask (Fig.~\ref{fig:production-studies}a).\\[3pt]
Exported geometry & Reconstruct faces from the SVG actually written to disk. &
Tests the exported partition, neck connectivity and planar interlock.\\
\bottomrule
\end{tabular}\end{table}\par
\textbf{Reading the rod measurements.} For the 14-piece Wave board
(seed~42), the closest eligible segments are 3.050\,mm apart on the accepted
rod centerlines snapped to a 0.01\,mm grid. The last retained collision sweep
recorded 2.723\,mm before a pair's repulsion correction. Both measurements
come from the solver. The SVG audit independently reconstructs the cut.
\clearpage
\section{Eight cuts on a common board}
\label{sec:export-comparison}
Formverk generated each style on the same $200\times160$\,mm rectangle,
with a target of 14 pieces and seed~42. Phase fields ran for 240 CPU steps
on a 0.8\,mm grid. Rods ran for 45 growth steps, with a fixed border for
Maze. Coral and Matrix used their default arc and mesh settings. Every run
used a 3.0\,mm minimum neck setting and a 0.15\,mm kerf.\par
\begin{figure}[!ht]\centering
\plate{cut-comparison}{1}
\caption{Interior cuts and perimeters read from the eight exported SVGs at
a common scale. Every style produced 14 recovered pieces and passed the
SVG audit for this seed.}
\label{fig:export-comparison}\end{figure}
\input{cut-comparison-table.tex}
\begin{table}[!ht]\centering\small
\caption{Measurements from independent SVG readback. Cut length includes the
perimeter and interior seams. A pass means that all mandatory geometric checks passed.}
\label{tab:cut-exports}
\begin{tabular}{@{}lrrrl@{}}\toprule
\textbf{Style} & \textbf{Pieces} & \textbf{Cut length (mm)} &
\textbf{Lock margin (mm)} & \textbf{Audit}\\\midrule
\cutcomparisonrows
\bottomrule\end{tabular}
\end{table}\par
For this seed, Dendrite, Amoeba and Linear differed in total cut length by
only 8.45\,mm across roughly 2.4\,m of cuts. Their planar lock margins
ranged from 0.891\,mm (Linear) to 3.141\,mm (Amoeba). The audit tests
360 translation directions and searches up to 12\,mm. Below that ceiling,
the search stops within a bracket of 0.0625\,mm.
\par\smallskip
\textbf{Reproduce the plate.} Run
\href{../../local/technical/compare_cuts.py}{\texttt{local/technical/compare\_cuts.py}}
from the repository. The accompanying
\href{../../local/technical/cut-comparison.csv}{\texttt{cut-comparison.csv}}
records the measurements and SHA-256 of each cut file.
\clearpage
\section{Discussion of adjacency and interlock}
\label{sec:interpretation}
Before any seam branches or bends, a seed construction can explain which
regions share a boundary. Consider seeds $\mathbf{s}_i$ in a planar board
$\Omega$. The ordinary Euclidean Voronoi cell of seed $i$ is
\begin{equation}
V_i=\{\mathbf{x}\in\Omega:\ \|\mathbf{x}-\mathbf{s}_i\|
\leq \|\mathbf{x}-\mathbf{s}_j\| \text{ for every }j\}.
\label{eq:voronoi}
\end{equation}
Along an interior edge shared by $V_A$ and $V_B$, the two distances are equal.
The corresponding line is the perpendicular bisector of segment $AB$.
A third seed limits how much of that line belongs to the shared edge.\par
\begin{figure}[!ht]\centering
\plate{bisector}{.33}
\caption{A point $P$ on the boundary has $PA=PB$. The dotted segments
are equal lengths. Moving either seed changes the bisector.}
\label{fig:bisect}\end{figure}\par
This construction explains adjacency and supplies a baseline for the diagrams.
The Dendrite paper uses reaction--diffusion initialization, and Wave begins
with an anisotropic layout. In either case, the initial territory and the
finished interlock answer different questions. The first determines which
pieces meet. Subsequent growth or optimization determines how they hold together.
\smallskip
\begingroup\small
\bibliographystyle{unsrt}
\bibliography{references}
\endgroup
\end{document}
